% ============================ 参考文献 ============================
\begin{thebibliography}{99}
\setlength{\itemsep}{1pt}
\small

\bibitem{grosch1994} Grosch W，Determination of Potent Odourants in Foods by Aroma
Extract Dilution Analysis (AEDA) and Calculation of Odour Activity Values (OAVs)，
Flavour and Fragrance Journal，9(4)：147-158，1994。

\bibitem{grosch1993} Grosch W，Detection of potent odorants in foods by aroma extract
dilution analysis，Trends in Food Science \& Technology，4(3)：68-73，1993。

\bibitem{caruana1997} Caruana R，Multitask Learning，Machine Learning，28(1)：41-75，1997。

\bibitem{chen2016} Chen T, Guestrin C，XGBoost: A Scalable Tree Boosting System，
Proceedings of the 22nd ACM SIGKDD International Conference on Knowledge Discovery and
Data Mining：785-794，2016。

\bibitem{breiman2001} Breiman L，Random Forests，Machine Learning，45(1)：5-32，2001。

\bibitem{wold2001} Wold S, Sjöström M, Eriksson L，PLS-regression: a basic tool of
chemometrics，Chemometrics and Intelligent Laboratory Systems，58(2)：109-130，2001。

\bibitem{lundberg2017} Lundberg S M, Lee S-I，A Unified Approach to Interpreting Model
Predictions，Advances in Neural Information Processing Systems 30：4765-4774，2017。

\bibitem{langfelder2008} Langfelder P, Horvath S，WGCNA: an R package for weighted
correlation network analysis，BMC Bioinformatics，9(1)：559，2008。

\bibitem{meinshausen2010} Meinshausen N, Bühlmann P，Stability selection，Journal of the
Royal Statistical Society: Series B (Statistical Methodology)，72(4)：417-473，2010。

\bibitem{kingma2014} Kingma D P, Welling M，Auto-Encoding Variational Bayes，
Proceedings of the 2nd International Conference on Learning Representations，2014。

\bibitem{dinh2017} Dinh L, Sohl-Dickstein J, Bengio S，Density estimation using Real NVP，
Proceedings of the 5th International Conference on Learning Representations，2017。

\bibitem{peyre2019} Peyré G, Cuturi M，Computational Optimal Transport with Applications
to Data Sciences，Foundations and Trends in Machine Learning，11(5-6)：355-607，2019。

\bibitem{horn1965} Horn J L，A rationale and test for the number of factors in factor
analysis，Psychometrika，30(2)：179-185，1965。
\bibitem{stone2004} Stone H, Sidel J L，Sensory Evaluation Practices，Amsterdam：
Elsevier Academic Press，3-23，2004。

\bibitem{kishimoto2018} Kishimoto T, Noba S, Yako N, Kobayashi M, Watanabe T，
Simulation of Pilsner-type beer aroma using 76 odor-active compounds，Journal of
Bioscience and Bioengineering，126(3)：330-338，2018。

\bibitem{yao2022} 姚孟琦，郭泽峰，郭世鑫，等，不同类型酵母对精酿啤酒化学和感官特性的影响，
中国酿造，41(1)：133-137，2022。

\bibitem{hao2007} 郝荣华，蔡宇杰，郑飞云，等，基于径向基函数神经网络探索风味物质对啤酒风味的
预测，中国酿造，26(9)：15-17，2007。

\bibitem{liu2019} 刘佳，黄淑霞，余俊红，等，基于电子舌技术的啤酒口感评价及其滋味信息与
化学成分的相关性研究，食品与发酵工业，45(2)：196-201，2019。

\bibitem{yang2023} Yang H, Wang Y, Zhao J, Li P, Li L, Wang F，A machine learning method
for juice human sensory hedonic prediction using electronic sensory features，Current
Research in Food Science，7：100576，2023。

\bibitem{feng2018} Feng Y, Cai Y, Fu X, Zheng L, Xiao Z, Zhao M，Comparison of
aroma-active compounds in broiler broth and native chicken broth by aroma extract
dilution analysis (AEDA), odor activity value (OAV) and omission experiment，
Food Chemistry，265：274-280，2018。

\end{thebibliography}
